\chapter{Spectral Yöntemler}
\label{ch:spectral}

Bu bölümde, periyodik sınır koşullarında PDE çözümünün en güçlü yöntemlerinden biri olan \textbf{spectral yöntemler}i inceliyoruz.
Temel fikir basit: fonksiyonu Fourier uzayina taşi, türev almak yerine çarpım yap, sonra geri dön.
Bu yaklaşim, sonlu fark (finite difference) yöntemlerine kiyasla \emph{üstel yakinsama} (exponential convergence) sağlar.


% ============================================================
\section{Fourier Serileri}
\label{sec:fourier_serileri}

\subsection{Gerçel Form}

Periyodik bir $f(x)$ fonksiyonunu ($f(x+L) = f(x)$) trigonometrik fonksiyonlarin toplami olarak yazabiliriz:

\begin{equation}
\label{eq:fourier_reel}
f(x) = \frac{a_0}{2} + \sum_{n=1}^{\infty} \left[ a_n \cos\!\left(\frac{2\pi n x}{L}\right) + b_n \sin\!\left(\frac{2\pi n x}{L}\right) \right]
\end{equation}

Burada $L$ periyot, $a_n$ ve $b_n$ Fourier katsayilari:
\begin{align}
a_n &= \frac{2}{L} \int_0^L f(x)\cos\!\left(\frac{2\pi n x}{L}\right) dx, \label{eq:an} \\
b_n &= \frac{2}{L} \int_0^L f(x)\sin\!\left(\frac{2\pi n x}{L}\right) dx. \label{eq:bn}
\end{align}

\begin{remark}
$a_0/2$ terimi fonksiyonun \textbf{ortalama değeri}dir: $a_0/2 = \frac{1}{L}\int_0^L f(x)\,dx$.
\end{remark}

\subsection{Kompleks (Exponential) Form}

Euler bağintisi $e^{i\theta} = \cos\theta + i\sin\theta$ kullanilarak, gerçel form çok daha zarif bir biçimde yazilabilir. $\cos$ ve $\sin$ terimlerini birleştirelim:
\begin{align}
a_n\cos(k_n x) + b_n\sin(k_n x) &= \frac{a_n - ib_n}{2}\,e^{ik_n x} + \frac{a_n + ib_n}{2}\,e^{-ik_n x}
\end{align}
Burada $k_n = 2\pi n / L$ dalga sayisidir. $\hat{c}_n \equiv (a_n - ib_n)/2$ tanimlanirsa:

\begin{equation}
\boxed{f(x) = \sum_{n=-\infty}^{\infty} \hat{c}_n \, e^{i k_n x}, \qquad k_n = \frac{2\pi n}{L}}
\label{eq:fourier_complex}
\end{equation}

\textbf{Fourier katsayilari} ortogonallik bağintisi $\frac{1}{L}\int_0^L e^{ik_n x}e^{-ik_m x}dx = \delta_{nm}$ kullanilarak:
\begin{equation}
\boxed{\hat{c}_n = \frac{1}{L}\int_0^{L} f(x)\,e^{-ik_n x}\,dx}
\label{eq:fourier_coeff}
\end{equation}

Bu kompleks form, türev alma işlemini dramatik biçimde basitleştirecek.


% ============================================================
\section{Ayrik Fourier Dönüşümü (DFT)}
\label{sec:dft}

Bilgisayarda sürekli Fourier integralini hesaplayamayız. Onun yerine fonksiyonu $N$ ayrik noktada örnekleriz.

\subsection{Grid Tanimlamasi}

$[0, L)$ araligini eşit araliklara bölelim:
\begin{equation}
x_j = j\,\Delta x, \qquad j = 0, 1, \ldots, N-1, \qquad \Delta x = \frac{L}{N}
\label{eq:grid}
\end{equation}

\begin{remark}
$N$ nokta ile $x_0 = 0$ ve $x_{N-1} = L - \Delta x$ olur. $x_N = L$ noktasi dahil \emph{edilmez} çünkü periyodik siniir koşulunda $f(0) = f(L)$ olduğundan bu nokta $x_0$ ile aynidir.
\end{remark}

\subsection{DFT Tanimi}

Sürekli Fourier integralinin ayrik karşiliği:

\begin{equation}
\boxed{\hat{f}_k = \sum_{j=0}^{N-1} f_j \, e^{-2\pi i j k / N}, \qquad k = 0, 1, \ldots, N-1}
\label{eq:dft}
\end{equation}

Ters dönüşüm (IDFT):
\begin{equation}
\boxed{f_j = \frac{1}{N}\sum_{k=0}^{N-1} \hat{f}_k \, e^{2\pi i j k / N}, \qquad j = 0, 1, \ldots, N-1}
\label{eq:idft}
\end{equation}

\begin{remark}
$1/N$ normalizasyonu IDFT'de. Bazi kaynaklarda (ve NumPy/PyTorch'ta) DFT'de $1/\sqrt{N}$ kullanilir. \textbf{Tutarlilik önemli}: DFT$\cdot$IDFT $= $ identity olmalı.
\end{remark}

\subsection{Dalga Sayilari ve Nyquist Frekansi}

$N$ noktayla temsil edilebilecek dalga sayilari sinirlidir:
\begin{equation}
k_n = \frac{2\pi n}{L}, \qquad n = -\frac{N}{2}, \ldots, -1, 0, 1, \ldots, \frac{N}{2} - 1
\label{eq:wavenumbers}
\end{equation}

\textbf{Nyquist frekansi}: temsil edilebilecek en yüksek frekans
\begin{equation}
k_{\max} = \frac{2\pi}{L}\cdot\frac{N}{2} = \frac{\pi}{\Delta x}
\label{eq:nyquist}
\end{equation}

Fiziksel yorum: en az 2 grid noktasi ile temsil edilebilecek en kisa dalga boyu $\lambda_{\min} = 2\Delta x$'tir.
Bu, \textbf{Shannon-Nyquist örnekleme teoremi}nin doğrudan sonuçudur.


% ============================================================
\section{FFT Algoritmasi}
\label{sec:fft}

\subsection{Hesaplama Karmaşıklığı}

DFT'yi (\ref{eq:dft}) doğrudan hesaplamak:
\begin{itemize}
\item Her $\hat{f}_k$ için $N$ çarpma + $N$ toplama $\Rightarrow$ toplam $O(N^2)$ işlem.
\item $N = 10^6$ için $10^{12}$ işlem --- pratik değil.
\end{itemize}

\subsection{Cooley-Tukey (1965): Böl ve Yönet}

FFT (Fast Fourier Transform), DFT'yi $O(N \log N)$'e düşüren algoritmadir.
Temel fikir: $N$-nokta DFT'yi iki $N/2$-nokta DFT'ye böl:

\begin{align}
\hat{f}_k &= \sum_{j=0}^{N-1} f_j\, e^{-2\pi i jk/N} \notag \\
&= \underbrace{\sum_{m=0}^{N/2-1} f_{2m}\, e^{-2\pi i mk/(N/2)}}_{\text{{çift indeksler}}: \; E_k}
\;+\;
e^{-2\pi ik/N}\;\underbrace{\sum_{m=0}^{N/2-1} f_{2m+1}\, e^{-2\pi i mk/(N/2)}}_{\text{tek indeksler}: \; O_k}
\label{eq:fft_decomp}
\end{align}

Yani:
\begin{equation}
\hat{f}_k = E_k + W^k \cdot O_k, \qquad W = e^{-2\pi i / N}\;\text{(twiddle factor)}
\label{eq:butterfly}
\end{equation}

Bu bölüme rekürsif olarak uygulanir:
\begin{equation}
\underbrace{N}_{\text{DFT}} \to 2\times\frac{N}{2} \to 4\times\frac{N}{4} \to \cdots \to N\times 1
\qquad \Rightarrow \quad \log_2 N \;\text{seviye} \times N \;\text{işlem/seviye} = O(N\log N)
\end{equation}

\begin{example}
Hiz karşilastirmasi ($N = 96 \times 160 \times 64 = 983{,}040$ --- bizim grid):
\begin{itemize}
\item Dogrudan DFT: $N^2 \approx 10^{12}$ işlem
\item FFT: $N\log_2 N \approx 2\times10^7$ işlem
\item \textbf{Hizlanma}: $\sim 50{,}000\times$
\end{itemize}
\end{example}

\subsection{3D FFT}

Üç boyutlu DFT, her boyutta ardisik 1D FFT'ler olarak hesaplanir (tensor çarpimi yapisi):
\begin{equation}
\hat{f}(k_x, k_y, k_z) = \text{FFT}_z\!\Big(\text{FFT}_y\!\big(\text{FFT}_x\!\big(f(x,y,z)\big)\big)\Big)
\end{equation}

PyTorch'ta tek satirla:
\begin{lstlisting}[caption={3D FFT PyTorch'ta},label={lst:fft3d}]
import torch.fft as fft
f_hat = fft.fftn(f, dim=(-3, -2, -1))  # 3 boyutta FFT
\end{lstlisting}


% ============================================================
\section{Spectral Türev Hesaplama}
\label{sec:spectral_deriv}

Spectral yöntemlerin as gizli silahi: \textbf{türev almak, Fourier uzayinda çarpmaya dönüşür}.

\begin{figure}[H]
\centering
\includegraphics[width=0.85\textwidth]{spectral_derivative.png}
\caption{Spectral türev hesaplama: fiziksel uzayda $f(x)$ ve $f'(x)$ (üst), Fourier uzayinda $|\hat{f}(k)|$ ve $|ik\hat{f}(k)|$ (alt). FFT $\to$ $ik$ ile çarpma $\to$ IFFT.}
\label{fig:spectral_derivative}
\end{figure}

\subsection{Temel Teorem: Birinci Türev}

\begin{theorem}[Spectral Türev]
\label{thm:spectral_deriv}
$f(x)$ periyodik ve pürüzsüz bir fonksiyon olsun. O zaman:
\begin{equation}
\boxed{f'(x) = \mathcal{F}^{-1}\!\big[\,ik\cdot\hat{f}(k)\,\big]}
\label{eq:spectral_deriv}
\end{equation}
Yani: $f'(x) = \text{IFFT}\big(ik \cdot \text{FFT}(f)\big)$
\end{theorem}

\begin{proof}[Türetme]
Fourier serisi temsili ile baslayalim:
\begin{equation}
f(x) = \sum_{n=-\infty}^{\infty} \hat{c}_n\, e^{ik_n x}
\end{equation}

Her iki tarafi $x$'e gore türetelim. Toplam ve türev yer değiştirebilir (üniform yakinsama):
\begin{align}
f'(x) &= \frac{d}{dx}\sum_{n=-\infty}^{\infty}\hat{c}_n\,e^{ik_n x} \notag \\
&= \sum_{n=-\infty}^{\infty}\hat{c}_n\cdot\frac{d}{dx}e^{ik_n x} \notag \\
&= \sum_{n=-\infty}^{\infty}\hat{c}_n\cdot (ik_n)\,e^{ik_n x} \notag \\
&= \sum_{n=-\infty}^{\infty}\underbrace{(ik_n\,\hat{c}_n)}_{\widehat{f'}(k_n)}\,e^{ik_n x}
\label{eq:spectral_deriv_proof}
\end{align}

Sonuçu okuyalim: $f'(x)$'in Fourier katsayilari $ik_n\,\hat{c}_n$'dir. Yani \textbf{türev almak} = \textbf{Fourier katsayisini $ik$ ile çarpmak}.
\end{proof}

\subsection{İkinci Türev}

Ayni mantigi iki kez uygulayelim:
\begin{align}
f''(x) &= \frac{d}{dx}f'(x) = \frac{d}{dx}\sum_n (ik_n)\hat{c}_n\,e^{ik_n x} \notag \\
&= \sum_n (ik_n)^2\,\hat{c}_n\,e^{ik_n x} = \sum_n (-k_n^2)\,\hat{c}_n\,e^{ik_n x}
\end{align}

\begin{equation}
\boxed{f''(x) = \mathcal{F}^{-1}\!\big[-k^2 \cdot \hat{f}(k)\big]}
\label{eq:spectral_2nd_deriv}
\end{equation}

Genel olarak $m$-inci türev:
\begin{equation}
\frac{d^m f}{dx^m} = \mathcal{F}^{-1}\!\big[(ik)^m \cdot \hat{f}(k)\big]
\label{eq:spectral_mth_deriv}
\end{equation}

\subsection{3D'de Gradient ve Laplacian}

Üç boyutlu $f(x,y,z)$ fonksiyonu için her yöndeki türev bağimsiz olarak hesaplanir. $\hat{f}(k_x,k_y,k_z)$ 3D Fourier dönüşümünü göstersin:

\textbf{Gradient:}
\begin{equation}
\pd{f}{x} = \mathcal{F}^{-1}\!\big[ik_x\cdot\hat{f}(k_x,k_y,k_z)\big], \qquad
\pd{f}{y} = \mathcal{F}^{-1}\!\big[ik_y\cdot\hat{f}\big], \qquad
\pd{f}{z} = \mathcal{F}^{-1}\!\big[ik_z\cdot\hat{f}\big]
\label{eq:3d_gradient}
\end{equation}

\textbf{Laplacian:}
\begin{align}
\laplacian f &= \pdd{f}{x} + \pdd{f}{y} + \pdd{f}{z} \notag \\
&= \mathcal{F}^{-1}\!\big[(-k_x^2)\hat{f}\big] + \mathcal{F}^{-1}\!\big[(-k_y^2)\hat{f}\big] + \mathcal{F}^{-1}\!\big[(-k_z^2)\hat{f}\big] \notag \\
&= \mathcal{F}^{-1}\!\big[-(k_x^2 + k_y^2 + k_z^2)\cdot\hat{f}\big]
\end{align}

\begin{equation}
\boxed{\laplacian f = \mathcal{F}^{-1}\!\Big[-|\vect{k}|^2 \cdot \hat{f}(\vect{k})\Big], \qquad |\vect{k}|^2 = k_x^2 + k_y^2 + k_z^2}
\label{eq:spectral_laplacian_3d}
\end{equation}

\subsection{Sayisal Örnek}

$L = 2\pi$ periyotlu $f(x) = \sin(2\pi x / L) = \sin(x)$ fonksiyonunun türevi $f'(x) = \cos(x)$ olduğunu biliyoruz. Spectral yöntemle doğrulayalım:

\begin{enumerate}
\item $f(x) = \sin(x)$ fonksiyonunun Fourier dönüşümü:
\begin{equation}
\sin(x) = \frac{e^{ix} - e^{-ix}}{2i} \qquad\Rightarrow\qquad \hat{c}_1 = \frac{1}{2i} = -\frac{i}{2},\quad \hat{c}_{-1} = \frac{i}{2}
\end{equation}
\item $ik$ ile çarpalim ($k_1 = 1$, $k_{-1} = -1$):
\begin{align}
ik_1 \cdot \hat{c}_1 &= i\cdot 1 \cdot\left(-\frac{i}{2}\right) = -\frac{i^2}{2} = \frac{1}{2} \\
ik_{-1}\cdot\hat{c}_{-1} &= i\cdot(-1)\cdot\frac{i}{2} = -\frac{i^2}{2} = \frac{1}{2}
\end{align}
\item Ters dönüşüm:
\begin{equation}
f'(x) = \frac{1}{2}\,e^{ix} + \frac{1}{2}\,e^{-ix} = \cos(x) \quad\checkmark
\end{equation}
\end{enumerate}


% ============================================================
\section{Anizotropik Domain}
\label{sec:anisotropic}

\subsection{Bizim Problem Geometrisi}

Projemizdeki periyodik kutu eşit kenarli değil:
\begin{equation}
\text{Domain:}\quad L_x = 6, \quad L_y = 10, \quad L_z = 4
\end{equation}
\begin{equation}
\text{Grid:}\quad N_x = 96, \quad N_y = 160, \quad N_z = 64
\end{equation}

\subsection{Yöne Bagli Dalga Sayilari}

Her yönün kendi periyodu ve grid çözünürlüğü oldugundan, dalga sayilari yönden yöne değişir:
\begin{equation}
k_x^{(n)} = \frac{2\pi n}{L_x}, \qquad
k_y^{(m)} = \frac{2\pi m}{L_y}, \qquad
k_z^{(l)} = \frac{2\pi l}{L_z}
\label{eq:aniso_wavenumbers}
\end{equation}

\begin{example}
Sayisal degerler:
\begin{itemize}
\item $x$ yonunde: $\Delta k_x = 2\pi/6 \approx 1.047$, \quad $k_{x,\max} = \pi N_x/L_x = \pi\cdot96/6 \approx 50.3$
\item $y$ yonunde: $\Delta k_y = 2\pi/10 \approx 0.628$, \quad $k_{y,\max} = \pi\cdot160/10 \approx 50.3$
\item $z$ yonunde: $\Delta k_z = 2\pi/4 \approx 1.571$, \quad $k_{z,\max} = \pi\cdot64/4 \approx 50.3$
\end{itemize}

Grid orani $N/L$ her yonde esit tutulmustur ($96/6 = 160/10 = 64/4 = 16$), bu da fiziksel cozunurlugu her yonde esitler.
\end{example}

\subsection{PyTorch Implementasyonu}

\texttt{innate.py} icerisindeki \texttt{SpectralOps3DAniso} sinifi tam olarak bu anizotropik yapıyı uygular:

\begin{lstlisting}[caption={SpectralOps3DAniso dalga sayisi hesabi},label={lst:aniso_wavenumbers}]
# Her yon icin ayri dalga sayisi dizisi
kx_1d = torch.fft.fftfreq(Nx, d=Lx/Nx) * 2*pi  # d = dx = Lx/Nx
ky_1d = torch.fft.fftfreq(Ny, d=Ly/Ny) * 2*pi
kz_1d = torch.fft.fftfreq(Nz, d=Lz/Nz) * 2*pi

# 3D meshgrid
kx, ky, kz = torch.meshgrid(kx_1d, ky_1d, kz_1d, indexing='ij')
\end{lstlisting}

\begin{remark}
PyTorch'un \texttt{fftfreq(N, d=dx)} fonksiyonu $[0, 1, 2, \ldots, N/2-1, -N/2, \ldots, -1] / (N \cdot dx)$ dondurur. $2\pi$ ile carpmak, fiziksel dalga sayisi $k_n = 2\pi n/L$'yi verir.
\end{remark}

\subsubsection{Nyquist Modu Sifirlama}

\texttt{SpectralOps3DAniso}'da ek bir incelik vardir: \textbf{Nyquist modunun sifirlanmasi}.
\begin{lstlisting}[caption={Nyquist modu sifirlama}]
if Nx % 2 == 0:
    kx_1d[Nx // 2] = 0.0  # Nyquist modu sifirla
\end{lstlisting}

$N$ cift iken $k_{N/2}$ modu belirsizdir: $e^{i\pi j}$ hem $\cos(\pi j)$ hem $-\cos(\pi j)$ olarak yorumlanabilir. Türev hesabinda $ik_{N/2}\hat{f}_{N/2}$ terimi yalancı (spurious) bir salınım yaratır. Sifirlamak bu ambiguity'yi yok eder ve türevin gerçel kalmasini garantiler.


% ============================================================
\section{De-aliasing}
\label{sec:dealiasing}

\subsection{Problem: Nonlineer Terimlerde Aliasing}

\begin{figure}[H]
\centering
\includegraphics[width=0.85\textwidth]{dealiasing.png}
\caption{De-aliasing: 2/3 kuralı. $|k| > N/3$ olan modlar sıfırlanır (kırmızı bölge). Bu, nonlineer çarpım sonrası aliasing hatasını önler.}
\label{fig:dealiasing}
\end{figure}

Navier-Stokes denklemindeki adveksiyon terimi $(\vect{u}\cdot\gradient)\vect{u}$ \textbf{nonlineerdir}: hız alanının kendisiyle çarpımı. Bu çarpım, Fourier uzayında \textbf{konvolüsyon}a karşılık gelir.

İki Fourier modunun carpimini düşünelim:
\begin{equation}
e^{ik_m x}\cdot e^{ik_n x} = e^{i(k_m + k_n)x}
\end{equation}

Eğer $|k_m| \leq k_{\max}$ ve $|k_n| \leq k_{\max}$ ise, çarpim $|k_m + k_n| \leq 2k_{\max}$ dalga sayisina sahip olabilir. Ama bizim gridimiz sadece $|k| \leq k_{\max} = N/2$'ye kadar temsil edebilir. $N/2 < |k_m + k_n| \leq N$ arasindaki modlar ne olur?

\textbf{Aliasing}: Temsil edilemeyen yüksek frekansli mod, düşük frekansli bir moda ``katlanir'':
\begin{equation}
e^{i(k_m + k_n)x_j} = e^{i(k_m + k_n - 2\pi N/L)x_j} \qquad\text{(grid üzerinde ayirt edilemez)}
\end{equation}

Bu, \textbf{sahte enerji} yaratir ve say\i{}sal çözümü bozar.

\subsection{2/3 Kurali --- Türetme}
\label{subsec:twothirds}

Aliasing'i önlemenin en yaygin yöntemi \textbf{2/3 kurali} (Orszag, 1971): hesaplamadan önce $|k| > N/3$ olan modlari sifirla.

\textbf{Neden tam olarak $N/3$?} Adim adim:

\begin{enumerate}
\item Carpimdan önce her iki çarpanin da $|k| \leq K$ modlari olsun.
\item Carpim sonucu $|k| \leq 2K$ modlari içerecektir.
\item Aliasing olmasin diye carpim sonucu gridin temsil sinirinda kalmali: $2K \leq k_{\max} = N/2$.
\item Dolayisyla $K \leq N/4$.
\end{enumerate}

Bir dakika --- bu $N/4$ verir, ama pratikte $N/3$ kullaniliyor. Fark şuradan gelir: padded convolution ile degil, \emph{truncated} Fourier uzayinda çalisiyoruz. Daha dikkatli analiz:

\begin{enumerate}
\item Fonksiyonlarimiz $[-N/2, N/2)$ arasinda modlara sahip.
\item $|k| > N/3$ modlarini sifirliyoruz, yani aktif modlar $|k| \leq N/3$.
\item İki aktif mod çarpimi: $|k_m + k_n| \leq 2N/3$.
\item $2N/3 < N/2$ olmadigi için aliasing olabilir, \emph{ama} $2N/3 < N$ oldugundan katlanma $N - 2N/3 = N/3$'ten büyük modlara gider --- bunlar zaten sifir!
\item Yani aliased enerji zaten sifirlanmis modlara duuser ve gerçek modlari \emph{kirletmez}.
\end{enumerate}

\begin{equation}
\boxed{\text{De-aliasing:}\quad \hat{f}_k = 0 \;\;\text{for}\;\; |k| > \frac{N}{3}}
\label{eq:dealias_rule}
\end{equation}

\begin{remark}
Etkin çözünürlük: $N = 96$ noktamiz varsa, sadece $N/3 = 32$ mod kullaniriz. Bu da cozunurlugun \%67'sini \emph{kullanmak}, \%33'ünü \emph{feda etmek} demektir. Pahali ama zorunlu.
\end{remark}

\subsection{3/2 Kurali (Padding Yöntemi)}

Alternatif yaklasim: çözünürlük kaybetmek yerine, çarpimi daha büyük bir gridde yap:

\begin{enumerate}
\item Fourier katsayilarini $3N/2$ boyutlu bir diziye zero-pad et.
\item $3N/2$ boyutlu IFFT ile fiziksel uzaya dön.
\item Fiziksel uzayda çarpımı yap ($f\cdot g$ noktasal).
\item $3N/2$ boyutlu FFT ile geri dön.
\item $|k| > N/2$ modlarini at, $N$ boyutuna geri dön.
\end{enumerate}

Bu yöntem hiç mod kaybettirmez ama \%50 daha fazla bellek ve hesaplama gerektirir. Pratikte iki yöntem birbirine denk sonuçlar verir.

\subsection{Kodda De-aliasing}

\texttt{SpectralOps3DAniso} icerisinde 2/3 kurali uygulamasi:
\begin{lstlisting}[caption={De-aliasing mask olusturma},label={lst:dealias}]
# Mod indeksleri uzerinden 2/3 kurali
mx = torch.fft.fftfreq(Nx, d=1.0) * Nx  # [-N/2, ..., N/2-1]
my = torch.fft.fftfreq(Ny, d=1.0) * Ny
mz = torch.fft.fftfreq(Nz, d=1.0) * Nz
Mx, My, Mz = torch.meshgrid(mx, my, mz, indexing='ij')

dealias_mask = (
    (torch.abs(Mx) < Nx // 3) &
    (torch.abs(My) < Ny // 3) &
    (torch.abs(Mz) < Nz // 3)
)
\end{lstlisting}


% ============================================================
\section{Poisson Denklemi Spectral Çözümü}
\label{sec:poisson}

\subsection{Problem}

Navier-Stokes çözümünde her zaman adiminda Poisson denklemi çözülmelidir:
\begin{equation}
\laplacian p = f \qquad\text{(verilen kaynak terimi $f$)}
\label{eq:poisson}
\end{equation}

Sonlu fark yöntemlerinde bu, büyük bir lineer sistem ($Ap = f$) demektir --- iteratif çözücü (CG, GMRES, multigrid) gerektirir. Spectral yöntemde ise \textbf{dogrudan çözüm} mümkündür.

\subsection{Türetme}

Her iki tarafa 3D FFT uygulayalim:
\begin{align}
\laplacian p = f \qquad &\xrightarrow{\text{FFT}} \qquad -(k_x^2 + k_y^2 + k_z^2)\,\hat{p} = \hat{f} \notag \\
&\xrightarrow{\text{çöz}} \qquad \hat{p}(\vect{k}) = \frac{-\hat{f}(\vect{k})}{|\vect{k}|^2}
\label{eq:poisson_spectral}
\end{align}

\begin{equation}
\boxed{p = \mathcal{F}^{-1}\!\left[\frac{-\hat{f}(\vect{k})}{k_x^2 + k_y^2 + k_z^2}\right]}
\label{eq:poisson_solution}
\end{equation}

\subsection{$k = 0$ Modu}

$\vect{k} = (0,0,0)$ için payda sifir olur. Bu mod fonksiyonun \textbf{uzaysal ortalamasi}dir:
\begin{equation}
\hat{p}(0,0,0) = \frac{1}{\text{Vol}}\int_\Omega p\,dV \equiv \langle p \rangle
\end{equation}

Periyodik siniir kosullarinda basinc $p$'ye bir sabit eklemek fizigi değiştirmez (çünkü sadece $\gradient p$ görülür). Bu yüzden:
\begin{equation}
\hat{p}(0,0,0) = 0 \qquad\text{(ortalama basinc sifir)}
\label{eq:mean_pressure}
\end{equation}

\subsection{Pressure Projection (Helmholtz Ayrismasi)}
\label{subsec:projection}

Navier-Stokes çözümünde kritik bir adim: hız alanını \textbf{divergence-free} (sıkıştırılamaz) yapmak.

\textbf{Helmholtz Ayrismasi Teoremi}: Her (pürüzsüz) vektör alan $\vect{v}$ iki bilesene ayrilabilir:
\begin{equation}
\vect{v} = \underbrace{\vect{u}}_{\text{solenoidal}\;(\divergence\vect{u}=0)} + \underbrace{\gradient\phi}_{\text{irrotasyonel}\;(\curl\gradient\phi=0)}
\label{eq:helmholtz}
\end{equation}

Projection (Chorin, 1968): divergence-free olmayan $\vect{v}$'den solenoidal $\vect{u}$'yu cikarma:
\begin{enumerate}
\item Divergence hesapla: $d = \divergence\vect{v}$
\item Poisson çöz: $\laplacian\phi = d$
\item Gradient çikar: $\vect{u} = \vect{v} - \gradient\phi$
\end{enumerate}

Spectral uzayda tamami tek satirda:
\begin{equation}
\hat{\vect{u}} = \hat{\vect{v}} - \frac{\vect{k}(\vect{k}\cdot\hat{\vect{v}})}{|\vect{k}|^2}
\label{eq:leray_projection}
\end{equation}

Bu \textbf{Leray projeksiyon operatörüdür}: $\mathbb{P} = \mathbb{I} - \vect{k}\vect{k}^T / |\vect{k}|^2$.

\begin{lstlisting}[caption={Spectral Poisson çözümü ve projection},label={lst:poisson}]
import torch.fft as fft

# Poisson cozumu: nabla^2 p = rhs
rhs_hat = fft.fftn(rhs)
k_sq = kx**2 + ky**2 + kz**2
k_sq_safe = k_sq.clone()
k_sq_safe[k_sq_safe == 0] = 1.0    # Divide-by-zero onle

p_hat = rhs_hat / (-k_sq_safe)
p_hat[..., 0, 0, 0] = 0.0          # Ortalama basinc = 0
p = fft.ifftn(p_hat).real

# Divergence-free projection
div_hat = 1j * (kx*u_hat + ky*v_hat + kz*w_hat)
phi_hat = div_hat / (-k_sq_safe)
phi_hat[..., 0, 0, 0] = 0.0
u_hat -= 1j * kx * phi_hat          # u = u - dp/dx
v_hat -= 1j * ky * phi_hat          # v = v - dp/dy
w_hat -= 1j * kz * phi_hat          # w = w - dp/dz
\end{lstlisting}


% ============================================================
\section{Spectral Yöntemlerin Avantaj ve Dezavantajlari}
\label{sec:spectral_tradeoffs}

\begin{table}[H]
\centering
\caption{Spectral vs. sonlu fark yöntemleri karsilastirmasi}
\label{tab:spectral_vs_fd}
\begin{tabular}{lcc}
\toprule
\textbf{Özellik} & \textbf{Spectral (FFT)} & \textbf{Sonlu Fark} \\
\midrule
Türev dogrulugu & Üstel (exponential) & Polinom ($O(\Delta x^p)$) \\
Hesaplama maliyeti & $O(N\log N)$ & $O(N)$ \\
Poisson çözümü & Dogrudan & İteratif \\
Siniir kosullari & Sadece periyodik & Her türlü \\
Karmasik geometri & Uygulanamaz & Uygulanabilir \\
Paralelleştirme & Zor (global FFT) & Kolay (lokal stencil) \\
De-aliasing gereksinimi & Evet & Hayir \\
\bottomrule
\end{tabular}
\end{table}

Bizim problemimiz tamamen periyodik (6 kenar da periyodik BC), bu da spectral yöntemi \textbf{ideal} kilar. Eger bir kenar bile periyodik olmasaydi (mesela duvar sinir kosulu), o yönde Chebyshev polinomlarina geçmek gerekirdi --- bu da ``pseudo-spectral'' yöntem olurdu.


% ============================================================
\section{Bu Bölümün Özeti}

\begin{enumerate}
\item Fourier serisi: periyodik fonksiyonları sinüs/kosinüs toplamı olarak yazar.
\item DFT: sürekli Fourier integralinin $N$-noktalı ayrık versiyonu.
\item FFT: DFT'yi $O(N^2) \to O(N\log N)$'e düşürür (Cooley-Tukey).
\item \textbf{Spectral türev}: Fourier uzayında $ik$ ile çarpmak $=$ fiziksel uzayda türev almak.
\item \textbf{Spectral Laplacian}: $-|\vect{k}|^2$ ile çarpmak.
\item \textbf{De-aliasing}: Nonlineer terimlerdeki sahte enerjiyi $2/3$ kuralıyla önle.
\item \textbf{Poisson denklemi}: $\hat{p} = -\hat{f}/|\vect{k}|^2$ --- doğrudan, iterasyonsuz çözüm.
\item \textbf{Projection}: Hız alanını divergence-free yapar (Helmholtz ayrışması).
\end{enumerate}

Bir sonraki bölümde, spectral yöntemlerle çözülen alanın fiziğini --- türbülans ve LES --- inceliyoruz.
